% Research CV. Publication dates, review statuses and authorship updated by the author on 5 November 2026.
% Layout follows the author-approved research CV; bibliography titles are published with author authorization.
\documentclass[11pt,letterpaper]{article}
\usepackage[utf8]{inputenc}
\usepackage[T1]{fontenc}
\usepackage[lf]{ebgaramond}
\usepackage[margin=0.68in,top=0.6in,bottom=0.62in]{geometry}
\usepackage{microtype}
\usepackage{enumitem}
\usepackage{titlesec}
\usepackage{fancyhdr}
\usepackage{xcolor}
\usepackage[hidelinks,pdfauthor={Shayan Hajhashemi},pdftitle={Research CV — Shayan Hajhashemi}]{hyperref}
\setlength{\parindent}{0pt}
\setlength{\parskip}{2pt}
\setlist[itemize]{leftmargin=1.15em,itemsep=1pt,topsep=2pt,parsep=0pt}
\titleformat{\section}{\large\bfseries}{}{0em}{}[\vspace{-4pt}\titlerule]
\titlespacing*{\section}{0pt}{9pt}{4pt}
\pagestyle{fancy}
\fancyhf{}
\renewcommand{\headrulewidth}{0pt}
\fancyfoot[L]{\scriptsize Shayan Hajhashemi · Updated 5 November 2026}
\fancyfoot[R]{\scriptsize\thepage}
\setlength{\footskip}{20pt}
\newlength{\projectheadingwidth}
\newcommand{\project}[2]{%
  \vspace{2pt}%
  \settowidth{\projectheadingwidth}{\textbf{#1}\quad\textit{#2}}%
  \ifdim\projectheadingwidth>\linewidth
    \textbf{#1}\par\nopagebreak[4]\textit{#2}\par\nopagebreak[4]%
  \else
    \textbf{#1}\hfill\textit{#2}\par\nopagebreak[4]%
  \fi}
\newcommand{\me}{\textbf{Hajhashemi, S.}}
\begin{document}
\begin{center}
{\LARGE\bfseries Shayan Hajhashemi}\\[3pt]
World models · Representation learning · Dynamical systems and control\\[3pt]
McGill University · Montréal, Canada\quad
\href{https://mila.quebec/en/directory/shayan-hajhashemi}{Mila profile}\\[2pt]
\small\href{https://shayshay42.github.io/}{Website}\quad
\href{https://github.com/shayshay42}{GitHub}\quad
\href{https://scholar.google.com/citations?user=uW3aYOcAAAAJ}{Google Scholar}\quad
\href{https://ca.linkedin.com/in/hshay}{LinkedIn}\\[2pt]
\footnotesize shayan [dot] hajhashemi [at] mail [dot] mcgill [dot] ca
\end{center}
\vspace{-11pt}

\section*{Research profile}
Scientific machine learning researcher working on representation learning and nonlinear dynamical systems. Research spans multimodal world models, learned inverse solvers, and constrained optimal control in differentiable simulations. Interested in action-conditioned representations and reliable world models for robot learning and planning.
Seeking research internships and full-time research roles.

\section*{Education}
\textbf{McGill University} — PhD, Quantitative Life Sciences\hfill Expected 2026\\
Supervisors: Morgan Craig and Amin Emad. Developing scientific machine learning tools in nonlinear dynamics and control with applications in pharmacology.\\
\textbf{McGill University} — BSc (Honours), Interdepartmental Immunology\hfill 2021\\
Honours research: Investigating Turing patterns in reaction--diffusion systems; supervisor: Anmar Khadra.

\section*{Selected research}
\project{Multimodal world models: alignment, prediction, and control}{AISTATS 2027 submission}
\begin{itemize}
\item Co-developed concepts for representation alignment and nonlinear model predictive control (NMPC) in multimodal world models; contributed to theoretical analysis.
\item Benchmarked our regularizer of learned representations to SINDy and time-series foundation models (TSFMs) on the same NMPC task for the Lorenz system.
\end{itemize}

\project{WeLD: label-free weak-form distillation for PDE inference}{Research draft}
\begin{itemize}
\item Coauthored a label-free method that learns weights over weak-form equations, distilling classical estimates from long observation windows into PDE parameter inference from short windows.
\item Analyzed teacher bias and instrumental-variable corrections; explored forecasting-based latent world models that learn parameter inference without a teacher.
\end{itemize}

\project{WAR: weak-form adaptive regularization for system identification}{AISTATS 2027 submission}
\begin{itemize}
\item Combined learned temporal-window selection and adaptive regularization with differentiable constrained weak-form solvers; used variable projection for nonlinear parameter inference from noisy observations.
\item Showcased superior inverse problem performance on held-out systems across six ODE families, evaluating parameter and trajectory recovery separately. Tested sensitivity to noise and sparse or irregular sampling.
\end{itemize}

\project{OIL: learning initializers through downstream optimization}{ICLR 2027 submission}
\begin{itemize}
\item Trained message-passing initializers through five differentiable trajectory-matching updates, with a three-term loss on initial parameters, refined parameters, and simulated trajectories.
\item Analyzed sensitivity geometry with Fisher information, investigated rollout failures and posterior diagnostics.
\end{itemize}

\project{DiffDose: differentiable optimal control for personalized medicine}{Preprint; under review at npj Systems Biology and Applications}
\begin{itemize}
\item Optimized bounded dose amplitudes and administration times through differentiable ODE and delay differential equation simulations. Compared automatic differentiation with analytic adjoints and finite differences for open-loop control.
\end{itemize}

\project{Latent-space dynamics and representation learning}{Preprint; under review at Nature Computational Science}
\begin{itemize}
\item Coauthored a variational framework that constrains learned cell-state representations of single-cell RNA-seq data to flows on an energy landscape, using potential-gradient neural ODEs for trajectory and cell-fate prediction.
\end{itemize}

\newpage
\section*{Publications and manuscripts}
\begin{enumerate}[leftmargin=1.2em,itemsep=7pt,topsep=2pt,parsep=0pt]
\item \me; Shum-Tim, C.; Craig, M.; Emad, A.
\textit{WAR: Weak-form Adaptive Regularization for Dynamical Systems Inverse Problems.}
Submitted 7 October 2026. Under review at \textit{AISTATS 2027}.
\item \me; Shum-Tim, C.; Prakash, S.; Emad, A.
\textit{Predators, Prey, and Sloppy Parameters: Fine-Tuning Amortized Inference for Generalized Lotka--Volterra.}
Submitted 18 September 2026. Under review at \textit{ICLR 2027}. \textbf{Co-first author.}
\item Ramezanian-Panahi, M.; \me; Rish, I.
\textit{Multimodal World Models: Alignment, Prediction, and Control.}
Submitted 7 October 2026. Under review at \textit{AISTATS 2027}.
\item \me; Emad, A.; Craig, M.
\textit{DiffDose: Differentiable Programming for Personalized Dose-Regimen Optimal Control.}
Submitted 13 September 2026. bioRxiv DOI: \href{https://doi.org/10.64898/2026.09.07.749974}{10.64898/\allowbreak 2026.09.07.749974}.
Under review at \textit{npj Systems Biology and Applications}.
\item Poursina, A.; \me; Mikaeili Namini, A.; Saberi, A.; Emad, A.; Najafabadi, H.~S.
\textit{A Latent Space Thermodynamic Model of Cell Differentiation.}
Submitted 6 March 2026. bioRxiv DOI: \href{https://doi.org/10.64898/2026.03.04.709512}{10.64898/\allowbreak 2026.03.04.709512}.
Under review at \textit{Nature Computational Science}.
\item \href{https://www.gene.com/scientists/our-scientists/vincent-lemaire}{Lemaire, V.}; \me; Rotte, A.; Seyed-Ahmadi, A.; Tripathi, S.; Chanu, P.; Wu, B.; Li, C.; Ramanujan, S.; Gadkar, K.; Jin, J.
\textit{Explainability analysis of QSP-derived virtual populations identifies robust pre-treatment biomarkers of immunotherapy.}
Accepted for publication in \textit{npj Precision Oncology}.
\end{enumerate}

\section*{Auxiliary publications}
\begin{enumerate}[start=7,leftmargin=1.2em,itemsep=7pt,topsep=2pt,parsep=0pt]
\item Brassard, A.; Dhoparee-Doomah, I.; \me; Bourdeau, F.; Giannias, B.; Spicer, J.; Ferri, L.~E.; Ma, K.; Cools-Lartigue, J.
\textit{Neutrophil Extracellular Trap Deposition Establishes an Immunosuppressive Environment in Non-metastatic Lymph Nodes and Facilitates Heterologous Metastatic Outgrowth.}
28 October 2026. bioRxiv, version 1. DOI: \href{https://doi.org/10.64898/2026.10.23.761387}{10.64898/\allowbreak 2026.10.23.761387}.

\item Dhoparee-Doomah, I.; Ma, K.; Brassard, A.; Leo, S.; \me; \ldots; Cools-Lartigue, J.
\textit{ROR$\gamma$ Inverse Agonism Redirects Tumour-Conditioned Emergency Myelopoiesis and Attenuates Established Tumour Progression.}
7 October 2026. bioRxiv, version 1. DOI: \href{https://doi.org/10.64898/2026.10.01.755741}{10.64898/\allowbreak 2026.10.01.755741}.
Under review at \textit{Cancer Discovery}.
\item Dhoparee-Doomah, I.; Li, Y.; Al Dow, M.; \me; \ldots; Cools-Lartigue, J.
\textit{A Measurable Systemic Immune Phenotype Links Circulating Myeloid Dysfunction to Tumour Spatial Organization in Gastroesophageal Adenocarcinoma.}
18 September 2026. bioRxiv, version 1. DOI: \href{https://doi.org/10.64898/2026.09.17.752366}{10.64898/\allowbreak 2026.09.17.752366}.
\item Roche, P.; \me; Akavia, U.~D. (2020).
\textit{Nanobody administration of Cas9 improves nuclear localization and HDR-enhanced synthetic lethality in A549 lung carcinoma.}
Withdrawn for intellectual property reasons. \href{https://www.jenthera.com/}{Jenthera Therapeutics}.
\end{enumerate}
Full author lists and patent: \href{https://shayshay42.github.io/publications.html}{shayshay42.github.io/publications.html}.

\newpage
\section*{Research experience}
\textbf{Genentech} — Modeling \& Simulations Intern\hfill 2024--2025\\
Supervisor: Vincent Lemaire; San Francisco, USA.
\begin{itemize}
\item Trained conditional generative models using flow matching and optimal transport for simulator-derived virtual populations; applied learned distributions to biomarker analysis.
\end{itemize}
\textbf{Earlier research:} Consultant Research Scientist, Jenthera Therapeutics (2020--2021); Research Associate, Goodman Cancer Research Center (2019--2020). Experimental nanobody research and computer vision.

\section*{Selected technical projects}
\textbf{Robot manipulation in simulation} — \href{https://github.com/shayshay42/robot}{Prototype for robot learning in simulation}, 2022--2023
\begin{itemize}
\item Built training and evaluation scripts for simulated manipulation with MuJoCo, robosuite, Gymnasium, and Stable-Baselines3; explored soft actor--critic and behavior cloning from demonstrations.
\end{itemize}
\textbf{Representation learning for dynamical systems and graphs} — McGill, 2021--2022
\begin{itemize}
\item Explored time-series embeddings using Transformers, LSTM autoencoders, Fourier and wavelet features for novelty search; studied relational graph convolutional embeddings and structural comparison of 3D RNA motifs.
\end{itemize}

\section*{Technical skills}
\textbf{Programming:} Python, Julia, MATLAB; PyTorch, JAX/diffrax, Pyro, Julia SciML.\\
\textbf{Methods:} Neural ODEs, graph neural networks, system identification, differentiable optimization, variational inference, flow matching, sensitivity analysis, model predictive control.\\
\textbf{Simulation projects:} MuJoCo, robosuite, Gymnasium, Stable-Baselines3; reinforcement and imitation learning experiments.

\section*{Presentations}
\begin{itemize}[leftmargin=1.15em,itemsep=6pt]
\item \me\ \href{https://meetings.siam.org/sess/dsp_talk.cfm?p=157552}{\textit{Real-MVP: Virtual Populations As Simulation-Based Conditional Neural Densities.}}
SIAM Conference on the Life Sciences, Cleveland, Ohio, 6 July 2026.
Invited oral presentation at an international conference. MS11: Recent Advances in Digital Twins and Trial Simulations to Guide Clinical Decisions.
\item Ramezanian-Panahi, M.; \me; Rish, I.
\href{https://www.main2026.org/abstracts}{\textit{Cross-modal control in the Virtual Brain.}}
Poster presentation (4D), \href{https://www.main2026.org/program}{MAIN 2026 — Montreal AI and Neuroscience}, HEC Montréal, Montréal, Canada, 30 October 2026.
\item Shum-Tim, C.; \me\ \textit{Weak-Form Distillation for Label-Free PDE Inverse Problems.}
\href{https://wmw-aspen.github.io/}{World Modeling for Physics Workshop}, Aspen Center for Physics, Aspen, Colorado, 28 February--5 March 2027.
\textbf{Accepted oral presentation; upcoming.}
\end{itemize}

\section*{Teaching and service}
\textbf{Teaching:} Scientific Machine Learning and ODEs in Julia, workshop instructor (2023); graduate nonlinear dynamics and computational methods, teaching assistant (2022).\\
\textbf{Service:} Julia conference reviewer (2024).

\section*{Selected awards}
NSERC CGS-D (2023--2026); FRQ-Santé doctoral scholarship (awarded 2023--2026); department three-minute thesis winner (2022); Dean's Multidisciplinary First Class Honours (2019--2021).
\end{document}
